\subsubsection{Subproceso 3.A: Verificaci\'on y Control de Reglas Acad\'emicas}

La revisi\'on documental y la entrevista semiestructurada identificaron los siguientes hallazgos sobre el funcionamiento efectivo del Subproceso 3.A:

\textbf{Actividades observadas y flujo efectivo:}
\begin{itemize}
    \item Alto grado de intervenci\'on manual en los casos de estudiantes convalidados o reincorporados (no siempre detectados autom\'aticamente por el SIGA en funci\'on de sus reglas).
\end{itemize}

\textbf{Tiempos de espera identificados:}
\begin{itemize}
    \item Existen cuellos de botella en la emisi\'on de dict\'amenes de rectificaci\'on (hasta 48-72 horas). Esto se debe al alto volumen de solicitudes (principalmente cruces de horarios).
\end{itemize}

\textbf{Problemas identificados:}
\begin{itemize}
    \item \textbf{Saturaci\'on en la direcci\'on de escuela por solicitudes repetitivas:} Estas solicitudes no est\'an prevenidas en el portal DIRCE, lo que provoca demoras en la aprobaci\'on de la Pre-Ficha y hace que se trasladen los desfases al subproceso 3.B.
\end{itemize}